% =========================================================
% Term Equality Rules
% =========================================================

\subsection{Deduction Rules: Term Equality}

% ---------------------------------------------------------
% TOOLKIT
% ---------------------------------------------------------

% ---------------------------------------------------------
% Term Equality Introduction (Reflexivity)
% ---------------------------------------------------------
\begin{definitionbox}{Definition (Term Equality Introduction --- Reflexivity)}
\begin{definition}[Term Equality Introduction --- Reflexivity]\label{def:eq-refl}
For any term $t$, one may infer $t = t$.
\end{definition}
\end{definitionbox}
\begin{remark*}[Standard quantified statement]
For every admissible object satisfying the hypotheses of the statement, the
assertion holds in the formal language under discussion.
\end{remark*}

\begin{remark*}[Interpretation]
The statement records the named logical notion or result together with its
intended syntactic or semantic role.
\end{remark*}


\begin{remark*}[Exposition]
Term reflexivity requires no premises and holds for all terms under all
interpretations. It is the local deduction rule that starts term-equality
reasoning.
\end{remark*}

\begin{dependencies}
\begin{itemize}
  \item \hyperref[def:term]{Term}
  \item \hyperref[ax:axiomatic-equality-reflexivity]{Reflexivity of Equality}
\end{itemize}
\end{dependencies}

% ---------------------------------------------------------
% Term Equality Elimination
% ---------------------------------------------------------
\begin{definitionbox}{Definition (Term Equality Elimination --- Substitution of Equals)}
\begin{definition}[Term Equality Elimination --- Substitution of Equals]\label{def:eq-elim}
From $t_1 = t_2$ and $\varphi[t_1/x]$, one may infer $\varphi[t_2/x]$.
\end{definition}
\end{definitionbox}
\begin{remark*}[Standard quantified statement]
For every admissible object satisfying the hypotheses of the statement, the
assertion holds in the formal language under discussion.
\end{remark*}

\begin{remark*}[Interpretation]
The statement records the named logical notion or result together with its
intended syntactic or semantic role.
\end{remark*}


\begin{remark*}[Exposition]
Term equality elimination permits replacing a term by an equal term in any
formula position, provided the substitution is capture-avoiding. The fuller
identity principle is developed separately; here the rule is recorded as part
of the first-order deduction toolkit.
\end{remark*}

% ---------------------------------------------------------
% Derived Rules
% ---------------------------------------------------------

\begin{dependencies}
\begin{itemize}
  \item \hyperref[def:term]{Term}
  \item \hyperref[ax:axiomatic-equality-substitution]{Substitution of Identicals}
\end{itemize}
\end{dependencies}

\begin{theorembox}{Theorem (Symmetry of Term Equality)}
\begin{theorem}[Symmetry of Term Equality]\label{thm:eq-sym}
From $t_1 = t_2$, one may infer $t_2 = t_1$.
\hyperref[prf:eq-sym]{Go to proof.}
\end{theorem}
\end{theorembox}
\begin{remark*}[Standard quantified statement]
$\forall t_1, t_2:\; t_1=t_2 \to t_2=t_1$.
\end{remark*}

\begin{remark*}[Predicate reading]
\[
  \forall t_1, t_2:\; t_1=t_2 \to t_2=t_1
\]
\end{remark*}



\begin{remark*}[Interpretation]
The predicate form says that term equality is symmetric: any ordered pair of
equal terms may be read in either order.

Symmetry licenses reversing an established term equality without adding new
mathematical information.
\end{remark*}

\begin{dependencies}
\begin{itemize}
  \item \hyperref[def:term]{Term}
  \item \hyperref[thm:axiomatic-equality-symmetry]{Symmetry of Equality}
\end{itemize}
\end{dependencies}

\begin{theorembox}{Theorem (Transitivity of Term Equality)}
\begin{theorem}[Transitivity of Term Equality]\label{thm:eq-trans}
From $t_1 = t_2$ and $t_2 = t_3$, one may infer $t_1 = t_3$.
\hyperref[prf:eq-trans]{Go to proof.}
\end{theorem}
\end{theorembox}
\begin{remark*}[Standard quantified statement]
$\forall t_1, t_2, t_3:\; t_1=t_2 \wedge t_2=t_3 \to t_1=t_3$.
\end{remark*}

\begin{remark*}[Predicate reading]
\[
  \forall t_1, t_2, t_3:\; t_1=t_2 \wedge t_2=t_3 \to t_1=t_3
\]
\end{remark*}



\begin{remark*}[Interpretation]
The predicate form says that term equality is transitive: two consecutive
term-equality links compose into one term-equality conclusion.

Transitivity is the rule that lets a proof chain substitutions through
intermediate terms.
\end{remark*}

\begin{dependencies}
\begin{itemize}
  \item \hyperref[def:term]{Term}
  \item \hyperref[def:eq-refl]{Term Equality Introduction}
  \item \hyperref[def:eq-elim]{Term Equality Elimination}
  \item \hyperref[thm:axiomatic-equality-transitivity]{Transitivity of Equality}
\end{itemize}
\end{dependencies}

\begin{theorembox}{Theorem (Function-Term Substitution under Term Equality)}
\begin{theorem}[Function-Term Substitution under Term Equality]\label{thm:eq-term}
If $t_1 = t_2$, then for any function symbol $f$,
$f(\dots,t_1,\dots) = f(\dots,t_2,\dots)$.
\hyperref[prf:eq-term]{Go to proof.}
\end{theorem}
\LeanFormalizes{thm:eq-term}{lra-lean}{LRA.VolumeI.Identity.Axioms}{SubstitutionPreservesFunctions}{pending}
\end{theorembox}
\begin{remark*}[Standard quantified statement]
$\forall t_1, t_2:\; t_1=t_2 \to \forall$ function symbol $f$: $f(\dots,t_1,\dots) = f(\dots,t_2,\dots)$.
\end{remark*}

\begin{remark*}[Predicate reading]
\[
  \forall t_1, t_2:\; t_1=t_2 \to \forall
\]
\end{remark*}



\begin{remark*}[Interpretation]
Equal input terms may be substituted inside a function-term expression without
changing the denoted output term.
\end{remark*}



\begin{dependencies}
\begin{itemize}
  \item \hyperref[def:eq-elim]{Term Equality Elimination}
  \item \hyperref[thm:axiomatic-equality-function-congruence]{Function Congruence}
\end{itemize}
\end{dependencies}

\begin{theorembox}{Theorem (Predicate Substitution under Term Equality)}
\begin{theorem}[Predicate Substitution under Term Equality]\label{thm:eq-pred}
If $t_1 = t_2$ and $P$ is an $n$-ary predicate symbol, then
$P(\dots,t_1,\dots) \;\Leftrightarrow\; P(\dots,t_2,\dots)$.
\hyperref[prf:eq-pred]{Go to proof.}
\end{theorem}
\LeanFormalizes{thm:eq-pred}{lra-lean}{LRA.VolumeI.Identity.Axioms}{SubstitutionPreservesPredicates}{pending}
\end{theorembox}
\begin{remark*}[Standard quantified statement]
$\forall t_1, t_2:\; t_1=t_2 \to \forall$ predicate symbol $P$: $P(\dots,t_1,\dots) \leftrightarrow P(\dots,t_2,\dots)$.
\end{remark*}

\begin{remark*}[Predicate reading]
\[
  \forall t_1, t_2:\; t_1=t_2 \to \forall
\]
\end{remark*}



\begin{remark*}[Interpretation]
Equal input terms may be substituted inside an atomic predicate without
changing the truth value of the predicate.
\end{remark*}



\begin{remark*}[Exposition]
These derived rules collectively express that function and predicate symbols
respect term equality: equal inputs produce equal outputs for function terms
and equivalent propositions for predicate formulas.
\end{remark*}

\begin{remark*}[Exposition]
When substituting equal terms, the substitution must be made consistently.
Replacing a term by an equal term in one occurrence but not another in the same
formula can change the intended statement.
\end{remark*}


\begin{dependencies}
\begin{itemize}
  \item \hyperref[def:eq-elim]{Term Equality Elimination}
  \item \hyperref[thm:axiomatic-equality-relation-congruence]{Relation Congruence}
\end{itemize}
\end{dependencies}
